\chapter{Final Considerations}
\label{chap:final_considerations}

This work addressed the problem of automatic scoring of ISSF paper shooting targets
using low-cost computer vision techniques. The motivation for the study arises from the need for
accessible scoring systems in amateur and training environments, where existing electronic
solutions are often too expensive or require specialized infrastructure.

An analysis of the physical characteristics and scoring rules of ISSF paper targets showed that
decimal scoring imposes high accuracy requirements, particularly for 10\,m air rifle
targets. A simple bounded-error model was introduced to illustrate how measurement accuracy
directly affects the probability of correct scoring, highlighting that even small localization
errors can result in score deviations. This analysis established the fundamental constraints that
any automatic scoring system must satisfy, independently of the specific algorithms employed.

The state of the art in computer-vision-based target scoring was reviewed, covering both
classical image-processing methods and recent deep learning approaches. Classical methods
rely on explicit geometric modeling and can achieve high precision under controlled conditions,
but they are sensitive to illumination changes, background clutter and overlapping
bullet holes, and often require careful parameter tuning. Deep learning-based approaches are
more robust and have low inference times, but the precision of their outputs is 
typically insufficient for high-accuracy ISSF
scoring when used alone. The literature review indicates that neither class of methods
alone fully addresses the combined requirements of robustness and precision.

Based on these observations, a hybrid system architecture was proposed. The system combines
robust detection mechanisms, potentially based on CNNs, with
geometric refinement stages designed to recover accurate bullet-hole and target center
coordinates. A modular image-processing pipeline was outlined, allowing different strategies
to be evaluated for target localization, bullet-hole detection, overlap handling and
geometric refinement.

The proposed architecture defines a clear structure for the dissertation phase, during which
the individual modules will be implemented, evaluated and integrated into a complete scoring
system. Future work will focus on selecting and refining for each module,
assessing their performance under realistic operating conditions, and
quantifying the scoring accuracy using both synthetic and real target data.

Overall, this work establishes a solid foundation for the development of a low-cost,
computer-vision-based automatic scoring system for ISSF paper targets, and clearly defines
the technical challenges and methodological directions to be addressed in the subsequent
implementation and evaluation stages.
